\section{Related Work}
\label{sec:related}

\paragraph{Tool use and trajectory evaluation.}
ReAct, Toolformer, and PAL integrate reasoning with tool calls or computation \citep{yao2023react,schick2023toolformer,gao2023pal}. ToolBench, API-Bank, and Gorilla study tool selection and invocation \citep{qin2023toolllm,li2023apibank,patil2024gorilla}; AgentBench and WebArena evaluate interactive agents \citep{liu2023agentbench,zhou2023webarena}. DAComp covers data engineering and analysis, including evaluator agreement and ranking stability \citep{lei2025dacomp}.

TRACE accumulates tool observations in an evidence bank and uses LLM evaluators to assess efficiency, hallucination, and adaptivity without requiring a single reference trajectory \citep{kim2025trace}. We share its evidence-centered, trajectory-level perspective. ToxicBench distinguishes support from observed evidence from correctness against unchanged source data: a conclusion can follow a corrupted observation yet be wrong. Our interventions track subsequent checking, poisoned-answer adoption, and recovery, complementing TRACE's trajectory-quality assessment.

\paragraph{Errors and recovery.}
\emph{Tools Fail} studies silent errors and recovery \citep{sun2024toolsfail}; ToolBench-X tests recoverable tool-environment hazards and compares diagnostic hints with additional inference \citep{tian2026toolbenchx}. ToolRobustBench aligns perturbations with tool interfaces, user intent, returned observations, and runtime feedback, using deterministic, cascade-aware failure diagnosis \citep{zheng2026toolrobustbench}. We share its controlled observation perturbations and process diagnostics. Its recovery setting is within a single tool-calling episode; ToxicBench follows later data-agent checks and final-answer adoption under one-shot and repeated poisoning, comparing ordinary retries, verification instructions, and later evidence conditions.

PALADIN and Expected Recovery Regret study recovery from tool failures and execution noise \citep{vuddanti2025paladin,vuddanti2026err}. AgentAtlas, structural testing, and AgentProp-Bench examine trajectory quality, component behavior, and error propagation \citep{mazaheri2026agentatlas,kohl2026structural,gurram2026agentprop}.

\paragraph{Poisoned instructions and evidence.}
Prompt injection embeds hostile instructions in user input or retrieved content \citep{perez2022ignore,greshake2023indirect}. MCP-ITP and TRUSTDESC address tool-metadata attacks and trustworthy descriptions \citep{li2026mcpitp,ye2026trustdesc}; AgentDojo, AgentCanary, and AgentSentry study agent attacks and defenses \citep{debenedetti2024agentdojo,li2026agentcanary,zhang2026agentsentry}. PoisonedRAG, SafeRAG, PoisonArena, and CPA-RAG study retrieval poisoning \citep{zou2024poisonedrag,liang2025saferag,chen2025poisonarena,li2025cparag}, while AgentDrift examines manipulated financial tool data \citep{wu2026agentdrift}. ToxicBench targets numerical and semantic corruption of returned data-tool observations (Table~\ref{tab:closest-work}).

\paragraph{Data quality and evidence checks.}
Data-validation systems check schemas and data quality \citep{schelter2018automating,breck2019data}; RADAR evaluates agents under tabular data-quality problems \citep{gu2025radar}. SDAG, RAGRank, and RAGShield address document attention, source credibility, and numerical-claim verification \citep{dekel2026sdag,jia2025ragrank,patil2026ragshield}. ToxicBench retains source tables as checking paths for altered tool responses.

\paragraph{Feedback and verification.}
Reflexion, Self-Refine, and MetaRAG use feedback or self-monitoring to revise answers \citep{shinn2023reflexion,madaan2023selfrefine,zhou2024metarag}. Co-Sight checks conflicts between structured facts, while Proof-of-Use trains agents to ground answers in retrieved evidence \citep{zhang2025cosight,ma2025pou}. ToolCritic provides feedback on specific tool-use errors \citep{hamad2025toolcritic}. Work on tool-use calibration and MICE links confidence estimates to tool decisions \citep{xuan2026confidence,subramani2025mice}; MIRROR evaluates whether models use self-knowledge to guide decisions \citep{wang2026mirror}. Our verification experiments compare ordinary second passes with prompts that ask for explicit checks, using the same route and token limits.
